% riemann-conjecture main.tex
% Research program: advancing the proportion of zeta zeros on the critical line
% Baseline: Claude/Anthropic (2026-08-10) — 2/3 unconditional

\documentclass[11pt,a4paper]{article}

\usepackage[utf8]{inputenc}
\usepackage[T1]{fontenc}
\usepackage{amsmath,amssymb,amsthm}
\usepackage{mathtools}
\usepackage{geometry}
\usepackage{hyperref}
\usepackage{cleveref}
\usepackage{booktabs}
\usepackage{tikz}
\usepackage{enumitem}
\usepackage[backend=biber,style=numeric]{biblatex}
\usepackage{microtype}

\geometry{margin=2.5cm}

% Theorem environments
\newtheorem{theorem}{Theorem}[section]
\newtheorem{lemma}[theorem]{Lemma}
\newtheorem{proposition}[theorem]{Proposition}
\newtheorem{corollary}[theorem]{Corollary}
\theoremstyle{definition}
\newtheorem{definition}[theorem]{Definition}
\theoremstyle{remark}
\newtheorem*{remark}{Remark}

% Notation
\newcommand{\zeta}{\zeta}
\newcommand{\R}{\mathbb{R}}
\newcommand{\C}{\mathbb{C}}
\newcommand{\N}{\mathbb{N}}
\newcommand{\Z}{\mathbb{Z}}
\newcommand{\Q}{\mathbb{Q}}
\newcommand{\F}{\mathbb{F}}
\newcommand{\re}{\operatorname{Re}}
\newcommand{\im}{\operatorname{Im}}
\newcommand{\tr}{\operatorname{tr}}
\newcommand{\norm}[1]{\left\|{#1}\right\|}
\newcommand{\abs}[1]{\left|{#1}\right|}
\newcommand{\Nzero}{N_0^*}
\newcommand{\Ntotal}{N}

\title{Advancing the Proportion of Riemann Zeta Zeros on the Critical Line: \\
A Multi-Model Research Program}

\author{Elcano Research Program\footnote{Automated research program using multi-model adversarial analysis. \\
Models: GLM-5.2 (orchestrator), Kimi K2.6, DeepSeek V4 Flash, Qwen 3.6, Cogito 2.1. \\
Repo: \texttt{github.com/elCanosail/riemann-conjecture} (to be made public).}}

\date{August 2026}

\begin{document}

\maketitle

\begin{abstract}
We report on a systematic research program to advance the unconditional lower bound on the proportion of non-trivial zeros of the Riemann zeta function $\zeta(s)$ that lie on the critical line $\re(s) = 1/2$. Our starting point is the result of Claude (Anthropic, 2026-08-10) that unconditionally at least $2/3$ of zeros lie on the critical line (optimised: $0.6725$ via the Montgomery--Taylor kernel), with a method ceiling of $0.68185$ for bandwidth-$1$ certificates. We survey all known approaches, verify the framework computationally, and conduct an adversarial multi-model analysis of potential improvements. We identify the key barriers and the most promising directions for future work.
\end{abstract}

\tableofcontents

\section{Introduction}

The Riemann Hypothesis (RH) asserts that all non-trivial zeros of $\zeta(s)$ lie on the critical line $\re(s) = 1/2$. While RH remains open, significant progress has been made on the weaker question: \emph{what proportion of zeros lie on the line?}

\subsection{Historical record}

\begin{center}
\begin{tabular}{lll}
\toprule
\textbf{Result} & \textbf{Constant} & \textbf{Method} \\
\midrule
Levinson (1974) & $\geq 1/3$ & Mollifier \\
Conrey (1989) & $\geq 5/12 \approx 0.4167$ & Mollifier (longer) \\
PRZZ (2020) & $> 5/12$ & Mollifier (refined) \\
\textbf{Claude (2026)} & $\geq \mathbf{2/3}$ \textbf{(0.6725 opt.)} & \textbf{Pair-corr. + linear algebra} \\
\bottomrule
\end{tabular}
\end{center}

\subsection{Claude's result}

On 2026-08-10, Anthropic announced that Claude proved unconditionally:
\begin{itemize}
    \item At least $2/3$ of non-trivial zeros lie on the critical line
    \item At least $2/3$ are simple and on the line
    \item At least $5/6$ are distinct
\end{itemize}
With the Montgomery--Taylor optimised kernel, these improve to $0.6725$, $0.6725$, and $0.83625$ respectively. The proof was formalised in Lean 4 against Mathlib's definition of $\zeta$.

\section{The framework}

[To be written: detailed exposition of Claude's method]

\section{Computational verification}

[To be written: numerical verification of zeta zeros and Gram matrix]

\section{Adversarial analysis}

\subsection{Can the rank-trace inequality be strengthened?}

[To be written: analysis of additional spectral data]

\subsection{Is the 0.68185 ceiling breakable?}

[To be written: analysis of certificate architectures]

\subsection{Can SDP help unconditionally?}

[To be written: semidefinite programming analysis]

\section{The frontier}

[To be written: synthesis of directions]

\section{Conclusion}

[To be written]

\printbibliography

\end{document}